\documentclass[10pt,a4paper]{amsart}
\usepackage[margin=19mm]{geometry}
\usepackage{amsmath,amssymb,mathtools}
\usepackage[scr=rsfs,cal=euler]{mathalfa}
\usepackage{xfrac}
\usepackage[unicode,hidelinks,bookmarks=false]{hyperref}
\newcommand{\ie}{i.e.\ }
\newcommand{\cf}{cf.\ }
\newcommand{\resp}{respectively\ }
\newcommand{\Subsection}{\subsection}
\providecommand{\nobreakdash}{\nobreak\hbox{-}}
\setlength{\emergencystretch}{2em}
\multlinegap0pt
\usepackage{amssymb}
\usepackage{enumerate}
\usepackage{booktabs}
\usepackage{xcolor}
\newtheorem{theorem}{Theorem}
\newtheorem{lemma}{Lemma}
\newtheorem{proposition}{Proposition}
\newcommand{\cut}{\backslash\backslash}
\title{A polynomial upper bound on Reidemeister moves for each link type}
\author{Marc Lackenby}
\date{Source: arXiv:2602.09923v1 (CC BY 4.0)}
\begin{document}
\maketitle

\noindent\emph{Source and licence.} A polynomial upper bound on Reidemeister moves for each link type. Version arXiv:2602.09923v1 (CC BY 4.0), distributed by arXiv under the Creative Commons Attribution 4.0 International licence, http://creativecommons.org/licenses/by/4.0/, as recorded in the arXiv licence metadata for this exact version and shown on the abstract page https://arxiv.org/abs/2602.09923v1. Full licence notice: \texttt{licenses/arxiv-2602.09923-CC-BY-4.0.txt}.\par\smallskip

\noindent\textit{Editorial scope.} Counted unit: the article's proposition Prop:AnnuliWeaklyFundamental (source0.tex lines 1072--1086). The article's complete original proof of it is delivered with it (source0.tex lines 1088--1166, one \texttt{proof} environment of 79 lines). No mathematics is written, repaired or re-derived: the statement and the proof are the article's own lines, in the article's own notation, and no retained line is substituted or re-flowed.  Retained uncounted as the source-local context the proof needs: lines 652--672; lines 1028--1041.  Nothing was omitted from any retained span, so the package carries no omission record.  No retained line contains a citation, so the wrapper carries no bibliography and no bibliography entry.  No retained line contains a figure environment, an image-inclusion command or drawing code, so no image is delivered and the package declares no asset dependency.  Editorial formatting only, in addition: an AMS \texttt{amsart} preamble that restates the article's own declarations and macros used by the retained lines, namely the environments \texttt{theorem}, \texttt{lemma}, \texttt{proposition} on separate counters, together with the standard packages those lines need.  The retained environments print in this document as Theorem 1, Lemma 1, Proposition 1 in order of appearance, whereas the article prints the same items with the numbering of its own full document.  The complete native source of the article as distributed in the arXiv e-print for this exact version is bundled unchanged as \texttt{sources/194356/source0.tex}.
\begin{theorem}
\label{Thm:SummandsEssentialEtc}
Let $(M,P)$ be a compact orientable irreducible boundary-irreducible 3-manifold with boundary pattern.
Let $\mathcal{T}$ be a triangulation of $M$ in which $P$ is simplicial. 
Let $S$ be a properly embedded, incompressible, boundary-incompressible, normal surface that has least possible weight,
among all normal surfaces that are pattern-isotopic to it. Assume either that $S$ is connected or that no component of $S$ is a sphere, clean disc, projective plane, boundary-parallel torus, or clean inessential annulus.
Suppose that $S$ can be expressed as a normal sum $S_1 + S_2$,
where $S_1$ and $S_2$ are in general position and non-empty, and the summation is in reduced form. Then the following hold.
\begin{enumerate}
\item No patch is a disc disjoint from $\partial M$ or a disc that intersects $\partial M$ in a single clean arc.
\item Every patch is incompressible and boundary-incompressible (with respect to the pattern $P$).
\item The trace surfaces are incompressible in $M \cut S$.
\item $S_1$ and $S_2$ are incompressible and boundary-incompressible (with respect to the pattern $P$).
\item No component of $S_1$ or $S_2$ is a sphere, clean disc or projective plane. 
\item If $S$ is orientable, then no component of $S_1$ or $S_2$ is a clean annulus that is parallel to an almost clean annulus in $\partial M$.
\item If $S$ is orientable and has least weight up to strong equivalence, then no component of $S_1$ or $S_2$ can be parallel to a toral component $T$ of $\partial M$
such that $|S \cap T| \leq 1$.
\item If $S$ is orientable, then no component of $S_1$ or $S_2$ can be parallel to a toral component $T$ of $\partial M$
disjoint from $P$.
\end{enumerate}
\end{theorem}

\begin{lemma}
\label{Lem:AnnuliInSFS}
Let $M$ be a Seifert fibre space with non-empty boundary, that embeds in the 3-sphere,
other than a solid torus or $S^1 \times S^1 \times I$.
\begin{enumerate}
\item If $\partial M$ is a single torus, then there is an essential annulus $A$ properly
embedded in $M$, such that the exterior of $A$ is two solid tori.
\item If $\partial M$ is more than one torus, let $T_1$ and $T_2$ be distinct
boundary components. Then there is an essential annulus $A$ properly embedded
in $M$, with one boundary component on $T_1$ and one boundary component on $T_2$.
\end{enumerate}
In both cases, $A$ is unique, up to a homeomorphism of $M$ that is
isotopic to the identity on $\partial M$.
\end{lemma}

\begin{proposition}
\label{Prop:AnnuliWeaklyFundamental}
Let $M$ and $A$ be as in Lemma \ref{Lem:AnnuliInSFS}. Let $P$ be a boundary pattern, such that for
each component $T$ of $\partial M$, one of the following holds:
\begin{enumerate}
\item $T$ is disjoint from $P$;
\item $T \cap P$ is a disjoint union of essential simple closed curves;
\item $T \cut P$ is discs.
\end{enumerate}
Suppose that among the annuli in Lemma \ref{Lem:AnnuliInSFS}, $A$ intersects $P$ as few times as possible.
When $A$ is disjoint from $P$ and $\partial M$ is connected, we also require that both components of
$\partial A$ lie in the same component of $\partial M \cut P$. Pick a triangulation of $M$ in which
$P$ is simplicial. Then $A$ is strongly equivalent to a
weakly fundamental surface.
\end{proposition}

\begin{proof}
Since $A$ is essential, it can be isotoped to a normal surface. We pick a 
representative for $A$ that has minimal weight, up to strong equivalence. 
Suppose that $A$ is not weakly fundamental.

Suppose first that $A$ can be expressed as a sum
$S_1 + S_2$ of normal surfaces, where $S_2$ is a sphere, projective plane, torus, clean disc or clean annulus. We may assume
that the sum is in reduced form.
By Theorem \ref{Thm:SummandsEssentialEtc}, these surfaces are both incompressible,
and neither has a sphere, projective plane, clean disc or clean inessential annulus component. Hence, $\chi(S_1) = \chi(S_2) = 0$.

Consider first the case when $S_2$, say, is a torus. Then, $\partial A = \partial S_1$.
So either some component of $S_1$ is an annulus or two components of $S_1$ are
M\"obius bands. But the latter case cannot arise, because each M\"obius band would have to be vertical and so would contain
a singular fibre of order 2, whereas $M$ contains at most two singular fibres and these have coprime orders.
So, $S_1$ has an essential annulus component. It cannot be boundary-parallel by our assumption on $P$.
Since it has the same boundary as $A$, it
is strongly equivalent to $A$. But $S_1$ has smaller weight than $A$, which is a contradiction.
So we may assume that $A$ has no torus summand.

Suppose now that $S_2$ is a clean essential annulus and that $S_1$ has an annulus component. 
Consider first the case where $M$ has base space a disc and two exceptional fibres.
Then, by our minimality assumption about $|A \cap P|$, $A$ must be disjoint from $P$.
So, by assumption, both components of $\partial A$ lie in the
same component of $\partial M \cut P$. Hence, $\partial S_1$ and
$\partial S_2$ also lie in this component. Hence, $S_2$ is strongly
equivalent to $A$. But it has smaller weight than $A$, which is a contradiction.
So suppose that $M$ has more than one boundary component. Then by assumption, the components of $\partial A$
lie on distinct tori $T_1$ and $T_2$. Now $\partial S_1 \cup \partial S_2$
lies on the same boundary components as $\partial A$. Since $S_2$ is clean,
the torus or tori $T$ containing $\partial S_2$ have a clean fibre. By our minimality 
assumption on $|A \cap P|$, $A \cap T$ is disjoint from $P$. Hence, $A \cap T$ lies
in the same component of $T \cut P$ as $S_1 \cap T$ and $S_2 \cap T$.
If $S_2$ has one boundary component
on $T_1$ and one boundary component on $T_2$, then we may replace $A$ by this surface.
So suppose that both boundary components of $S_2$ lie on some boundary
component $T$ of $\partial M$, then $\partial S_1 \cap (\partial M - T) = \partial A \cap (\partial M - T)$,
which is single curve. This curve cannot be the boundary of a M\"obius band component of $S_1$,
since then every other component of $S_1$ is a torus or an annulus with both boundary curves on $T$,
which implies that $[A \cap T] = [S_1 \cap T] + [S_2 \cap T] =0 \in H_1(T; \mathbb{Z}_2)$,
but this is not the case.
So we may replace $A$ by the annular component of $S_1$ containing $\partial A \cap (\partial M - T)$
and reduce its weight, which is a contradiction.

Suppose now that $S_2$ is a clean essential annulus and that $S_1$ has no annulus or torus components.
Then $S_1$ is a collection of M\"obius bands. But $M$ cannot support two or more disjoint M\"obius bands.
So $S_1$ is a single M\"obius band.
If $M$ has a single boundary component, then, as above, $S_2$ is strongly equivalent to $A$ and has smaller
weight, which is a contradiction. So, suppose that $M$ has more than one boundary component,
and that one component of $\partial A$ lies on $T_1$ and the other lies on $T_2$.
Then $\partial S_1$ is a single regular fibre on some component ($T_1$, say) of
$\partial M$. The curves $\partial S_2$ are isotopic to regular fibres on $\partial M$.
If both curves of $\partial S_2$ lie on $T_2$, then $[A \cap T_2] = [S_2 \cap T_2] = 0 \in H_1(\partial M; \mathbb{ Z}_2)$,
which is not the case. If one curve of $\partial S_2$ lies on $T_1$ and the other lies on $T_2$,
then $[A \cap T_1] = [\partial S_1 \cap T_1] + [\partial S_2 \cap T_1] = 0 \in H_1(\partial M; \mathbb{ Z}_2)$, which 
again is not the case.
Finally if both curves of $\partial S_2$ lie on $T_1$, then $A$ is disjoint from $T_2$, which is not the case.
In all cases, we obtain a contradiction.

Suppose now that $A = S_1 + S_2 + S_3$ where $S_2$ and $S_3$ are clean M\"obius bands. Suppose
also that $M$ has a single boundary component.
Then $S_2$ (say) is a vertical surface that projects to an arc in the base space running from
the boundary of the base space to a singular point of order 2. So $2 S_2$ is a clean essential annulus.
By assumption, $A$ is therefore clean and both its boundary components lie in the
same component of $\partial M \cut P$. Hence, the annuli $2S_2$ and $2S_3$
also have boundary lying in this component, and so are strongly equivalent to $A$.
Now $w(S_2 + S_3)$ is the mean of $w(2S_2)$ and $w(2S_3)$.
So one of $w(2S_2)$ and $w(2S_3)$ is at most $w(S_2 + S_3)$, the former say. But
then $w(2S_2) \leq w(S_2 + S_3) < w(A)$, contradicting the assumption that $A$ has least weight.

Suppose now that $A = S_1 + S_2 + S_3$ where $S_2$ and $S_3$ are clean M\"obius bands, and that
$M$ has at least two boundary components. Now $[\partial S_2]$ represents the homology class
of a fibre of $T_1$ or $T_2$. The same is true of $[\partial S_3]$. Hence, $S_2 + S_3$
cannot be a single M\"obius band and some tori. It also cannot have more than one
M\"obius band component, since $M$ does not support two disjoint M\"obius bands.
So, $S_2 + S_3$ contains a clean annulus component. We may therefore write $A = S_1' + S'_2$
where $S'_2$ is a clean annulus, and we have dealt with that case above.

\end{proof}

\end{document}
